\chapter{Methodologie}%
\label{ch:methodologie}

Om de onderzoeksvraag te kunnen beantwoorden, wordt een proof of concept ontwikkeld dat uit een e-mail een interpreteerbare risicoscore berekent. Dit hoofdstuk geeft een overzicht van de fasen waarin het onderzoek is verlopen. Elke fase, behalve de literatuurstudie, wordt in een eigen hoofdstuk uitgewerkt.

\begin{enumerate}
    \item \textbf{Requirements en ontwerpkeuzes (Hoofdstuk~\ref{ch:requirements}).} Op basis van de literatuurstudie worden de functionele en inhoudelijke vereisten voor de risicoscore vastgelegd. Hiermee wordt deelvraag 1 beantwoord. De literatuurstudie en de afbakening van de onderzoeksscope leveren daarnaast de theoretische onderbouwing voor deelvraag 2. De requirements vormen vervolgens het uitgangspunt voor de verdere afbakening en het ontwerp van het proof of concept.
    \item \textbf{Voorbereiding van de data (Hoofdstuk~\ref{ch:data}).} Er wordt een dataset samengesteld met zowel phishingmails als legitieme e-mails. Na de nodige voorbereiding wordt deze opgesplitst in een trainings- en testset, zodat het model onafhankelijk kan worden geëvalueerd.
    \item \textbf{Kenmerkextractie (Hoofdstuk~\ref{ch:kenmerkextractie}).} Vervolgens worden mechanismen ontwikkeld om relevante kenmerken automatisch uit e-mails te extraheren. Elke extractor levert niet alleen een kenmerkwaarde op, maar registreert ook het bijbehorende bewijs in de oorspronkelijke e-mail. Deze fase biedt antwoord op deelvraag 3.
    \item \textbf{Uitleg (Hoofdstuk~\ref{ch:uitleg}).} De geëxtraheerde kenmerken en hun bijdrage aan de risicoscore worden vertaald naar niet technische verklaringen. Hiervoor worden uitlegregels opgesteld die aangeven welke elementen uit een e-mail aan de uiteindelijke beoordeling bijdragen. Hiermee wordt ook deelvraag 4 van antwoord voorzien.
    \item \textbf{Evaluatie (Hoofdstuk~\ref{ch:evaluatie}).} Tot slot wordt nagegaan in welke mate het proof of concept aan de vooropgestelde succescriteria voldoet.
\end{enumerate}

Deze volgorde is bewust gekozen, omdat elke fase voortbouwt op de vorige. De requirements bepalen de ontwerpkeuzes, die op hun beurt bepalen welke data en kenmerken nodig zijn. De uitleg is gekoppeld aan de kenmerken, en de evaluatie toetst of de requirements zijn gehaald.

\chapter{Requirements en ontwerpkeuzes}%
\label{ch:requirements}

In dit hoofdstuk worden eerst de eisen vastgelegd waaraan de risicoscore moet voldoen. Daarna volgen de ontwerpkeuzes die uit die eisen en uit de literatuurstudie (Hoofdstuk~\ref{ch:stand-van-zaken}) voortvloeien, samen met de afbakening van het onderzoek. Hoofdstuk~\ref{ch:data} bouwt hierop voort door de dataset te beschrijven waarop het model wordt getraind.

\section{Requirements}%
\label{sec:requirements}

Deelvraag 1 vraagt aan welke eisen een risicoscore moet voldoen om gebruikers te ondersteunen bij het beoordelen van e-mails. Het effect op gebruikers valt buiten de scope van dit onderzoek. De eisen worden afgeleid uit sectie~\ref{subsec:uitleg-voor-gebruikers}.

\begin{description}
  \item[R1: weinig \emph{false positives}.] Gebruikers wennen aan foutieve waarschuwingen die ze herhaaldelijk zien \autocite{Anderson2015}. False positives moeten daarom zo veel mogelijk vermeden worden.
  \item[R2: aanwijzingen bij de bron.] Waarschuwingen werken beter dicht bij de verdachte link \autocite{Petelka2019}. De uitleg moet gebruikers dus kunnen wijzen op de plaats in de e-mail waar de aanwijzing staat.
  \item[R3: begrijpelijke taal.] Veel gebruikers begrijpen technische waarschuwingen niet \autocite{Egelman2008}. De uitleg moet dus in gewone taal en zonder jargon zijn geschreven.
  \item[R4: geen overdreven zekerheid.] Een overtuigende uitleg kan tot \emph{overreliance} leiden \autocite{Bucinca2021}. De uitleg mag dus geen overdreven zekerheid uitstralen en moet overeenkomen met de werkelijke beoordeling.
\end{description}

\section{Ontwerpkeuzes en scopeafbakening}%
\label{sec:afbakening-scope}

De oplossing kent aan een e-mail een risicoscore toe die aangeeft hoe sterk de aangeboden input op phishing wijst. De input volgt het formaat van RFC~5322 \autocite{Resnick2008}.

Gelet op de eis van uitlegbaarheid, wordt gekozen voor een interpreteerbaar model. Zoals besproken in sectie~\ref{subsec:machinelearning-detectie:logistische-regressie} is logistische regressie een geschikte kandidaat, aangezien elk kenmerk een gewicht krijgt dat rechtstreeks aangeeft hoe het de risicoscore beïnvloedt. Verder kan er voor elke beslissing vooraf een uitleg worden gedefinieerd die dan in het uiteindelijke antwoord kan worden gecapteerd.

De kenmerken waarop het model getraind zal worden, worden berekend uit de datasets die worden besproken in Hoofdstuk~\ref{ch:data}.

Wat buiten de scope van dit onderzoek valt en dus niet wordt opgenomen:
\begin{enumerate}
    \item \textbf{Bijlagen als detectiekenmerk.} Zoals aangehaald in sectie~\ref{subsec:bijlagen} blijken bijlagekenmerken volgens \textcite{Mendes2025} slechts een beperkte bijdrage te hebben aan het eindoordeel van het model. De geringe meerwaarde weegt niet op tegen de bijkomende complexiteit;
    \item \textbf{LLM's voor beoordeling en uitleg.} Op basis van de bevindingen van \textcite{Turpin2023} valt een LLM af als generator van uitleg, aangezien de uitleg waarheidsgetrouw moet zijn, en dat kan een LLM niet in elk geval garanderen. Ook \emph{indirect prompt injection} \autocite{Greshake2023} (sectie~\ref{subsec:machinelearning-detectie:nadelen}) vormt een risico voor de correcte werking van de oplossing. De inhoud van een phishingmail wordt immers door een aanvaller bepaald en kan dus ongewenste instructies voor de LLM bevatten;
    \item \textbf{Inhoudelijke (lexicale) kenmerken.} GenAI stelt foutloze en overtuigende teksten makkelijk op, waardoor taalkenmerken minder betrouwbaar worden dan structurele kenmerken. Verder zijn de beschikbare corpora overwegend Engelstalig \autocite{Pajola2025} (sectie~\ref{subsec:datasets}), wat maakt dat een model dat op woordgebruik wordt getraind niet voor een andere taal kan worden ingezet;
    \item \textbf{Stilometrische kenmerken.} Deze kenmerken vergelijken het schrijfgedrag in een bericht met de gebruikelijke schrijfstijl van de afzender. Dat vereist voldoende eerdere berichten van dezelfde afzender \autocite{Stringhini2014} (sectie~\ref{subsec:inhoud-taalgebruik:stilometrische-kenmerken}), terwijl het algoritme telkens één afzonderlijk bericht beoordeelt;
    \item \textbf{Authenticatiekenmerken.} De resultaten van SPF-, DKIM- en DMARC-controles (sectie~\ref{subsec:afzender-headers:authenticatiekenmerken}) worden niet als kenmerk opgenomen. De legitieme e-mails uit het SpamAssassin Public Corpus dateren van vóór de wijdverspreide invoering van deze authenticatiemechanismen en bevatten daardoor doorgaans geen \texttt{Authentication-Results}-header, terwijl de recente phishingmails uit de dataset die header vaak wel bevatten. Het model zou de aanwezigheid van die header dan als signaal voor phishing kunnen leren, terwijl ze enkel het tijdperk van de e-mail weerspiegelt. Moderne legitieme e-mails zouden daardoor ten onrechte als phishing worden aangemerkt, wat de \emph{false-positive-rate} (R1) zou verhogen.
\end{enumerate}
\chapter{Voorbereiding van de data}%
\label{ch:data}

Dit hoofdstuk beschrijft de dataset waarop het model wordt getraind en getest, en hoe die wordt voorbereid. De keuze van de dataset volgt uit de beperkingen van de bestaande corpora uit sectie~\ref{subsec:datasets} en uit de scopeafbakening in sectie~\ref{sec:afbakening-scope}. De voorbereide data vormen de invoer voor de kenmerkextractie in Hoofdstuk~\ref{ch:kenmerkextractie}.

\section{Gebruikte datasets}%
\label{sec:gebruikte-datasets}

De in sectie~\ref{subsec:datasets} besproken corpora zijn niet geschikt voor dit onderzoek, omdat ze niet de vereiste e-mailheaders bevatten. Daarom wordt geen gebruik gemaakt van MeAJOR \autocite{Mendes2025}, Enron \autocite{klimt2004} of Nazario\footfullcite{Nazario2025}.

De dataset waarop de oplossing wordt getraind, wordt opgebouwd uit de volgende twee datasets:
\begin{enumerate}
    \item \textbf{Phishing Pot.} Een actief onderhouden verzameling phishingmails \autocite{PhishingPot}, aangeleverd als \emph{.eml}-bestanden die rechtstreeks kunnen worden verwerkt;
    \item \textbf{SpamAssassin Public Corpus.} Een verzameling e-mailberichten uit de vroege jaren 2000, met zowel ongewenste als legitieme e-mails\footfullcite{SpamAssassinCorpus}. Voor dit onderzoek wordt uitsluitend gebruikgemaakt van de partities \emph{easy\_ham} en \emph{hard\_ham}, die beide legitieme e-mails bevatten. Deze partities worden eerst geconsolideerd, zoals beschreven in sectie~\ref{sec:consolidatie-spamassassin}.
\end{enumerate}

\section{Consolidatie van het SpamAssassin-corpus}%
\label{sec:consolidatie-spamassassin}

De geselecteerde partities van het SpamAssassin Public Corpus worden samengebracht in één dataset. Elk bericht bestaat uit een tekstbestand met de ruwe e-mail en een bijbehorend JSON-bestand met metadata, waaronder de MD5-checksum van het tekstbestand. Omdat voor dit onderzoek enkel de ruwe berichten nodig zijn, brengt een Bash-script de tekstbestanden samen in één map en verwijdert het de JSON-bestanden.

Tijdens deze consolidatie wordt voor elk bericht gecontroleerd of de MD5-checksum van het tekstbestand overeenkomt met de checksum in het bijbehorende JSON-bestand. Eén bericht uit \emph{*hard\_ham\_1*} vertoonde een afwijkende checksum en wordt daarom uitgesloten, aangezien de integriteit ervan niet kan worden aangetoond. De resulterende dataset bevat 4149 legitieme e-mails: 2500 uit \emph{easy\_ham\_1}, 1400 uit \emph{easy\_ham\_2} en 249 uit \emph{hard\_ham\_1}.

Omdat de identificatienummers enkel binnen een partitie uniek zijn, krijgt elk bestand de naam van zijn oorspronkelijke partitie als prefix, bijvoorbeeld \texttt{easy-ham-1\_00001.txt}.

Tot slot genereert het script een bestand \texttt{manifest.jsonl} met één JSON-object per bericht. Dit bevat onder meer de bestandsnaam, de partitie, het label (legitiem) en de checksum. Het manifest registreert zo de herkomst en integriteit van elk bericht en dient als centrale referentie voor de verdere voorbereiding van de dataset.


\section{Voorbereiding}%
\label{sec:voorbereiding}

Na de eventuele consolidatie van de bronbestanden wordt elke e-mail ingelezen en omgezet naar een verwerkbare vorm. Indien nodig worden headers met \emph{encoded-words} (RFC~2047 \autocite{Moore1996}) of een base64-gecodeerde body eerst gedecodeerd.

Vervolgens wordt de dataset opgesplitst in een trainingsset van 70\% en een testset van 30\%. Hierbij blijft de verhouding tussen phishingmails en legitieme e-mails in beide sets behouden. Een vaste seed zorgt ervoor dat de opsplitsing bij herhaaldelijke uitvoering reproduceerbaar is.

Uit elk bericht worden vervolgens de kenmerken berekend. Deze kenmerken, samen met de indeling van de berichten in trainings- en testset, worden opgeslagen. Zo hoeven de berichten niet bij elk experiment opnieuw te worden verwerkt, en worden alle modelconfiguraties op dezelfde gegevens getraind.

De parameters van het model, waaronder de beslissingsdrempel, worden afgesteld via \emph{cross-validation} op de trainingsset. Hierbij wordt de trainingsset in gelijke deelsets opgesplitst en wordt het model telkens getraind op alle deelsets behalve één, waarop het vervolgens wordt gevalideerd.

De drempel is de score vanaf waar een e-mail als phishing wordt aangeduid. Ze wordt zo gekozen dat de FPR in de validatiedeelsets het streefdoel van 3\% niet overschrijdt (sectie~\ref{sec:requirements}, R1). De testset wordt pas daarna gebruikt, uitsluitend om het eindresultaat te rapporteren.

\section{Beperkingen van de data}%
\label{sec:beperkingen-data}

Voor zover bekend is er geen recente, publiek beschikbare dataset van legitieme e-mails die de volledige, ongewijzigde headers bewaart. Het SpamAssassin Public Corpus is daarom de enige bruikbare bron van legitieme e-mails met headers, ook al dateert het van begin jaren 2000 en raadt het project zelf aan om de legitieme set niet meer te gebruiken wanneer die ouder is dan 18 maanden.\footfullcite{SpamAssassinCorpusPolicy}

Het risico bestaat dat het hedendaagse mailverkeer op punten afwijkt van deze verouderde data, wat de resultaten van het model kan beïnvloeden. Dit potentiële risico wordt meegenomen in de evaluatie van het model.
\chapter{Kenmerkextractie}%
\label{ch:kenmerkextractie}

In dit hoofdstuk wordt beschreven welke kandidaat-kenmerken uit de e-mails worden gehaald en hoe de extractors zijn opgebouwd. De kenmerken volgen uit de literatuurstudie (sectie~\ref{sec:kenmerken-detectie}) en uit de afbakening in sectie~\ref{sec:afbakening-scope}. De extractors werken op de data uit Hoofdstuk~\ref{ch:data}. Hun uitvoer vormt de basis voor de uitleg in Hoofdstuk~\ref{ch:uitleg}.

\section{Kandidaat-kenmerken}%
\label{sec:kandidaat-kenmerken}

Uit elk bericht worden de kenmerken afgeleid die in sectie~\ref{sec:afbakening-scope} binnen de scope vallen, voor zover ze rechtstreeks uit de datasets van Hoofdstuk~\ref{ch:data} komen of eruit af te leiden zijn. Het gaat om vier families die besproken werden in sectie~\ref{sec:kenmerken-detectie}.

De kenmerken worden eerst allemaal berekend als kandidaat-kenmerken. Welke ervan in het uiteindelijke model blijven, wordt bepaald aan de hand van L1-regularisatie, besproken in sectie~\ref{subsec:machinelearning-detectie}.

\paragraph{Opbouw van de extractors.}
Elke extractor levert twee elementen: een numerieke waarde voor het model en het concrete bewijs in de e-mail waarop die waarde berust. Zo kan de uitleg later controleerbaar naar de e-mail zelf verwijzen, wat requirement R2 (sectie~\ref{sec:requirements}) vertaalt.

\paragraph{Afzenderidentiteit.}
De kenmerken in Tabel~\ref{tab:kenmerken-afzender} volgen uit sectie~\ref{subsec:afzender-headers}. De controle op domeinen die lijken op legitieme domeinen en de controle op merken in de weergavenaam (A7) hangen af van een referentielijst.

\begin{table}[htbp]
    \centering
    \footnotesize
    \setlength{\tabcolsep}{4pt}
    \begin{tabular*}{\linewidth}{@{\extracolsep{\fill}} l p{0.7\linewidth} p{0.2\linewidth} @{}}
        \toprule
        \textbf{ID} & \textbf{Kenmerk} & \textbf{Waarde} \\
        \midrule
        \textbf{A1} & \texttt{From} wijkt af van \texttt{Reply-To} & boolean \\
        \textbf{A2} & \texttt{From} wijkt af van \texttt{Return-Path} & boolean \\
        \textbf{A3} & \texttt{Reply-To} wijkt af van \texttt{Return-Path} & boolean \\
        \textbf{A4} & Aantal verschillende domeinen over A1--A3 & int (1--3) \\
        \textbf{A5} & Weergavenaam bevat een adres of domein dat afwijkt van het echte adres & boolean \\
        \textbf{A6} & Weergavenaam deelt geen token met het adres & boolean \\
        \textbf{A7} & Merknaam in weergavenaam, maar niet in het domein & boolean \\
        \textbf{A8} & Tekens uit meerdere alfabetten in het \texttt{From}-domein & int \\
        \textbf{A9} & \texttt{From}-domein is geen geldig registreerbaar domein & boolean \\
        \textbf{A10} & Lengte van het \texttt{From}-domein & int \\
        \textbf{A11} & Cijfers in het \texttt{From}-domein & int \\
        \textbf{A12} & Aantal koppeltekens in \texttt{From}-domein & int \\
        \textbf{A13} & Labels in het \texttt{From}-domein & int \\
        \bottomrule
    \end{tabular*}
    \caption[Kandidaat-kenmerken afzenderidentiteit]{\label{tab:kenmerken-afzender}Kandidaat-kenmerken voor afzenderidentiteit.}
\end{table}

\paragraph{Transportkenmerken.}
De kenmerken in Tabel~\ref{tab:kenmerken-transport} worden eveneens besproken in sectie~\ref{subsec:afzender-headers:transport}.

\begin{table}[htbp]
    \centering
    \footnotesize
    \setlength{\tabcolsep}{4pt}
    \begin{tabular*}{\linewidth}{@{\extracolsep{\fill}} l p{0.7\linewidth} p{0.2\linewidth} @{}}
        \toprule
        \textbf{ID} & \textbf{Kenmerk} & \textbf{Waarde} \\
        \midrule
        \textbf{B1} & Domein van de \texttt{Message-ID} wijkt af van het \texttt{From}-domein & boolean \\
        \textbf{B2} & \texttt{Message-ID} is ongeldig opgebouwd & boolean \\
        \textbf{B3} & Domein van de \texttt{Message-ID} is geen geldig domein & boolean \\
        \textbf{B4} & Aantal \texttt{Received}-headers na opschoning & int \\
        \textbf{B5} & Verschil tussen \texttt{Date} en oudste \texttt{Received}-tijdstempel & float \\
        \textbf{B6} & \texttt{Received}-tijdstempels zijn niet monotoon & boolean \\
        \bottomrule
    \end{tabular*}
    \caption[Kandidaat-kenmerken transport]{\label{tab:kenmerken-transport}Kandidaat-kenmerken voor transport.}
\end{table}

Kenmerk B4 telt het aantal hops in de \texttt{Received}-keten. Kenmerk B5 meet het verschil tussen het \texttt{Date}-veld en het oudste \texttt{Received}-tijdstempel. Omdat \texttt{Date} door de afzender zelf wordt ingesteld en de \texttt{Received}-tijdstempels door de betrokken servers worden toegevoegd, kan een groot verschil wijzen op manipulatie. Kenmerk B6 controleert ten slotte of de tijdstempels van opeenvolgende hops chronologisch oplopen, aangezien een afwijkende volgorde kan wijzen op een toegevoegde of aangepaste \texttt{Received}-header.

\paragraph{Lexicale URL-kenmerken.}
De kenmerken in Tabel~\ref{tab:kenmerken-url} volgen uit sectie~\ref{subsec:urls-hyperlinks}.

\begin{table}[htbp]
    \centering
    \footnotesize
    \setlength{\tabcolsep}{4pt}
    \begin{tabular*}{\linewidth}{@{\extracolsep{\fill}} l p{0.7\linewidth} p{0.2\linewidth} @{}}
        \toprule
        \textbf{ID} & \textbf{Kenmerk} & \textbf{Waarde} \\
        \midrule
        \textbf{C1} & Aantal unieke URL's & int \\
        \textbf{C2} & Maximale URL-lengte & aantal tekens \\
        \textbf{C3} & Maximaal aantal subdomeinen & int \\
        \textbf{C4} & Maximaal aantal punten in de host & int \\
        \textbf{C5} & Maximaal aantal koppeltekens in de host & int \\
        \textbf{C6} & \texttt{@} in de URL & boolean \\
        \textbf{C7} & Verdachte woorden in host, pad of gebruikersinfo & int (woordenlijst) \\
        \textbf{C8} & Merknaam in host, pad of gebruikersinfo buiten het registreerbare domein & boolean (referentielijst) \\
        \textbf{C9} & Registreerbaar domein valt onder de private sectie van de Public Suffix List & boolean \\
        \textbf{C10} & Host is een IP-adres & boolean \\
        \textbf{C11} & Meerdere alfabetten in de link-host & boolean \\
        \bottomrule
    \end{tabular*}
    \caption[Kandidaat-kenmerken lexicale URL-kenmerken]{\label{tab:kenmerken-url}Kandidaat-kenmerken voor lexicale URL-kenmerken.}
\end{table}

C9 steunt op de private sectie van de Public Suffix List, die onder meer gratis hostingplatformen bevat die ook door aanvallers worden misbruikt.

\paragraph{Misleidende hyperlinks.}
De kenmerken in Tabel~\ref{tab:kenmerken-hyperlinks} volgen uit sectie~\ref{subsec:urls-hyperlinks:misleidende-hyperlinks}.
Deze kenmerken zijn enkel van toepassing op e-mails waarvan de body een HTML-deel bevat. Bij platte tekst is de weergegeven tekst tegelijk de bestemming, terwijl deze bij HTML van elkaar kunnen afwijken.

\begin{table}[htbp]
    \centering
    \footnotesize
    \setlength{\tabcolsep}{4pt}
    \begin{tabular*}{\linewidth}{@{\extracolsep{\fill}} l p{0.7\linewidth} p{0.2\linewidth} @{}}
        \toprule
        \textbf{ID} & \textbf{Kenmerk} & \textbf{Waarde} \\
        \midrule
        \textbf{D1} & Link mismatch: getoonde tekst komt niet overeen met achterliggende bestemming & boolean? \\
        \textbf{D2} & Geëncodeerde host in de \texttt{href} & boolean \\
        \textbf{D3} & Valse tooltip: \texttt{title}-attribuut bevat een URL met ander registreerbaar domein dan de \texttt{href} & boolean \\
        \bottomrule
    \end{tabular*}
    \caption[Kandidaat-kenmerken misleidende hyperlinks]{\label{tab:kenmerken-hyperlinks}Kandidaat-kenmerken voor misleidende hyperlinks.}
\end{table}

\section{Niet-bepaalbare kenmerken}%
\label{sec:niet-bepaalbare-kenmerken}

Sommige kandidaat-kenmerken zijn niet voor elke e-mail te bepalen, bijvoorbeeld omdat het bericht geen HTML-deel of geen \texttt{Reply-To} bevat. In dat geval wordt het kenmerk expliciet als ``niet van toepassing'' vastgelegd. Dit mag niet gelijkgesteld worden met een waarde van 0 of \emph{false}, want dit betekent dat een controle werd uitgevoerd en geen afwijking vond. Omdat logistische regressie geen ontbrekende waarden kan verwerken, worden de kenmerken gestandaardiseerd en wordt een niet-bepaalbaar kenmerk op de gemiddelde waarde van de trainingsdata gezet. Daardoor draagt het niets bij aan de risicoscore en wordt het niet als reden weergegeven.
\chapter{Uitleg}%
\label{ch:uitleg}

Dit hoofdstuk beschrijft hoe de risicoscore wordt omgezet in een uitleg voor de eindgebruiker. De uitleg bouwt voort op de kenmerken en het bewijs uit Hoofdstuk~\ref{ch:kenmerkextractie} en op de requirements uit sectie~\ref{sec:requirements}. In Hoofdstuk~\ref{ch:evaluatie} wordt nagegaan in hoeverre de uitleg aan de gestelde uitgangspunten voldoet.

\section{Uitgangspunten}%
\label{sec:uitgangspunten-uitleg}

De uitleg beantwoordt de vraag waarom het systeem een e-mail als risicovol beoordeelt en wordt in gewone taal aan de eindgebruiker voorgelegd. Zoals besproken in sectie~\ref{subsec:vormen-van-uitlegbaarheid} is hiervoor een lokale uitleg het meest geschikt, aangezien die de beoordeling van één individuele e-mail verduidelijkt \autocite{Guidotti2018}.

Bij het opstellen van de uitlegregels worden de volgende uitgangspunten gehanteerd, die requirement R2 tot en met R4 vertalen uit sectie~\ref{sec:requirements}:
\begin{enumerate}
    \item \textbf{Suggestief.} De verantwoordelijkheid voor de interpretatie blijft bij de eindgebruiker. De uitlegregels mogen dus geen definitief oordeel over de e-mail uitbrengen, maar dienen de gebruiker met voorzichtige formuleringen zoals ``kan wijzen op'' te informeren;
    \item \textbf{Waarheidsgetrouw.} Elke reden komt rechtstreeks uit de berekening van de score en berust op kenmerken die in het bericht daadwerkelijk worden vastgesteld. Kenmerken die niet of onvoldoende aanwezig zijn, worden niet als reden opgenomen;
    \item \textbf{Controleerbaar.} Elke reden verwijst naar concreet bewijs uit het bericht. Dat bewijs wordt samen met de kenmerkwaarde door de extractors bijgehouden (Hoofdstuk~\ref{ch:kenmerkextractie}), zodat de gebruiker elke reden zelf in de e-mail kan nagaan;
    \item \textbf{Toegankelijk.} De uitleg is geschreven in gewone taal en vermijdt jargon.
\end{enumerate}

\section{Kanttekening bij uitleg in deze bachelorproef}%
\label{sec:kanttekening-uitleg}

De uitleg is een essentieel onderdeel van het algoritme, want ondersteuning van de gebruiker vraagt niet enkel om een score, maar ook om de redenen erachter (sectie~\ref{sec:uitlegbaarheid-gebruikersondersteuning}). Dit onderzoek toont aan dat die redenen rechtstreeks uit de berekening van de score kunnen worden afgeleid, naar concreet bewijs in de e-mail verwijzen en in gewone taal kunnen worden weergegeven. De formulering wordt niet bij gebruikers getoetst en het effect op hun begrip, vertrouwen en beslissingen, inclusief \emph{overreliance} (sectie~\ref{subsec:uitleg-voor-gebruikers}), valt buiten de scope.
\chapter{Evaluatie}%
\label{ch:evaluatie}

Dit hoofdstuk beschrijft hoe het proof of concept wordt geëvalueerd. Eerst komt de evaluatie van de classificatieprestaties aan bod, daarna die van de uitleg. De evaluatie toetst de succescriteria uit de onderzoeksdoelstelling (sectie~\ref{sec:onderzoeksdoelstelling}) en de requirements uit sectie~\ref{sec:requirements}.

section{Classificatieprestaties}%
\label{sec:evaluatie-classificatie}

Om de classificatieprestaties van het algoritme te evalueren, wordt gebruikgemaakt van de testset die tijdens de voorbereiding van de data werd afgesplitst (Hoofdstuk~\ref{ch:data}). De resultaten van het model worden gerapporteerd aan de hand van de metrieken uit sectie~\ref{subsec:evaluatiemetrieken}, met de false-positive-rate (\emph{FPR}) als de belangrijkste, conform R1 uit sectie~\ref{sec:requirements}. Deze metriek geeft immers aan hoe vaak het algoritme een legitieme e-mail ten onrechte als phishing aanduidt.

Zoals besproken in sectie~\ref{subsec:uitleg-voor-gebruikers}, besteden gebruikers steeds minder aandacht aan waarschuwingen die ze herhaaldelijk te zien krijgen \autocite{Anderson2015}. Een hoge \emph{FPR} brengt dan ook het risico mee dat gebruikers aan de waarschuwingen van het algoritme wennen, waardoor dit onderzoek zijn doel mist.

\section{Kwaliteit van de uitleg}%
\label{sec:evaluatie-uitleg}

Naast de classificatieprestaties wordt nagegaan of de uitleg voldoet aan de uitgangspunten \emph{waarheidsgetrouw} en \emph{controleerbaar} uit sectie~\ref{sec:uitgangspunten-uitleg}. Daarmee worden requirement R2 en het waarheidsgetrouwe deel van requirement R4 uit sectie~\ref{sec:requirements} getoetst.

Om dit te rapporteren, wordt op een willekeurige steekproef van 50 e-mails uit de testset nagegaan of elk kenmerk dat bijdraagt tot de risicoscore effectief in de e-mail wordt vastgesteld en verwijst naar concreet bewijs in het oorspronkelijke bericht. De uitleg voldoet aan beide uitgangspunten wanneer dit voor alle bijdragende kenmerken in alle e-mails uit de steekproef het geval is.

De uitgangspunten \emph{suggestief} en \emph{toegankelijk} betreffen daarentegen de formulering van de uitleg en kunnen enkel via gebruikers worden getoetst. Dit valt buiten de scope van dit onderzoek, waardoor requirement R3 en het suggestieve deel van requirement R4 niet worden getoetst.